\section{Experiments}

\subsection{Environment Settings}

The experiments are conducted in a two-dimensional heterogeneous UAV cooperative pursuit environment with three pursuer UAVs and one maneuvering target. The pursuers have different maximum speeds, sensing ranges, and energy-related parameters. The target follows a mixed maneuvering policy, which combines escaping from the nearest pursuer and random turning. The final main evaluation uses the mixed target policy, target speed 0.75, communication radii $4,6,8,10$, 300 evaluation episodes per seed, and three random seeds.

\begin{table*}[t]
\centering
\caption{Training and evaluation settings used in the reported experiments.}
\label{tab:training_settings}
\begin{tabularx}{\textwidth}{lcl}
\toprule
Item & Value & Scope \\
\midrule
Parallel environments & 8 & All training runs \\
Rollout length & 128 steps & All training runs \\
PPO epochs & 4 & All training runs \\
Discount factor $\gamma$ & 0.99 & All training runs \\
GAE $\lambda$ & 0.95 & All training runs \\
Clip coefficient & 0.2 & All training runs \\
Entropy coefficient & 0.01 & All training runs \\
Value loss coefficient & 0.5 & All training runs \\
Max gradient norm & 0.5 & All training runs \\
Learning rate & $3\times10^{-4}$ & Default training; staged fine-tuning may use a lower value \\
Hidden dimension & 128 & Actor, critic, and graph encoders \\
Role embedding dimension & 8 & GAT-MAPPO and EA-RG-MAPPO-S \\
Graph minibatch size & 256 graphs & GAT-MAPPO and EA-RG-MAPPO-S \\
MAPPO minibatch size & 512 samples & MAPPO baseline \\
Target policy & mixed & Final evaluation and robustness tests \\
Target speed & 0.75 & Final main evaluation \\
Training communication radius & 8 & Fixed-radius stage \\
Staged fine-tuning radius & $U(4,10)$ & EA-RG-MAPPO-S random-radius stage \\
Final evaluation episodes & 300 per seed & Main table \\
Evaluation seeds & 3 & All reported mean$\pm$std tables \\
Python environment & cac / Python 3.8.20 & Reported runtime \\
PyTorch & 2.4.1+cu124 & Reported runtime \\
GPU & NVIDIA GTX 1650 Ti & Reported runtime \\
\bottomrule
\end{tabularx}
\end{table*}


Table~\ref{tab:training_settings} summarizes the training and evaluation settings. All methods use the centralized training and decentralized execution setting. MAPPO uses a shared actor and a centralized critic. GAT-MAPPO adds node-level graph attention on top of MAPPO. EA-RG-MAPPO-S further introduces role embeddings, relative edge features, and staged random-radius fine-tuning. During training, the policy first learns basic coordination under a fixed communication radius $R_c=8$. It is then fine-tuned from the fixed-radius checkpoint under randomly sampled radii $R_c\sim U(4,10)$.

\subsection{Compared Methods}

We compare MAPPO, GAT-MAPPO, and EA-RG-MAPPO-S. MAPPO is a multi-agent reinforcement learning baseline without an explicit graph structure. GAT-MAPPO adds standard graph attention to the MAPPO framework. EA-RG-MAPPO-S is the final proposed method with edge-aware role graph encoding and staged random-radius fine-tuning. For ablation analysis, RG-MAPPO and EA-RG-MAPPO are also included to analyze the effects of role graph modeling, edge-aware attention, and staged fine-tuning.

\subsection{Main Results}

\begin{table*}[t]
\centering
\caption{Final 300-episode evaluation under different communication radii. Results are reported as mean$\pm$std over three seeds.}
\label{tab:final_comm_300_results}
\begin{tabular}{lccccc}
\toprule
Method & Radius & Success $\uparrow$ & Collision $\downarrow$ & Timeout $\downarrow$ & Avg. steps $\downarrow$ \\
\midrule
MAPPO & 4 & 0.707$\pm$0.167 & 0.228$\pm$0.099 & 0.072$\pm$0.070 & 84.632$\pm$31.747 \\
MAPPO & 6 & 0.770$\pm$0.145 & 0.218$\pm$0.130 & 0.017$\pm$0.019 & 70.081$\pm$29.733 \\
MAPPO & 8 & 0.810$\pm$0.136 & 0.177$\pm$0.126 & 0.017$\pm$0.015 & 62.143$\pm$25.158 \\
MAPPO & 10 & 0.836$\pm$0.077 & 0.152$\pm$0.067 & 0.014$\pm$0.011 & 58.812$\pm$16.622 \\
\midrule
GAT-MAPPO & 4 & 0.811$\pm$0.026 & 0.136$\pm$0.020 & 0.060$\pm$0.040 & 72.822$\pm$14.039 \\
GAT-MAPPO & 6 & 0.838$\pm$0.051 & 0.126$\pm$0.034 & 0.039$\pm$0.025 & 67.784$\pm$18.906 \\
GAT-MAPPO & 8 & 0.784$\pm$0.054 & 0.179$\pm$0.042 & 0.038$\pm$0.042 & 71.418$\pm$18.202 \\
GAT-MAPPO & 10 & 0.802$\pm$0.037 & 0.156$\pm$0.024 & 0.043$\pm$0.034 & 71.886$\pm$12.111 \\
\midrule
EA-RG-MAPPO-S & 4 & 0.926$\pm$0.004 & 0.054$\pm$0.007 & 0.022$\pm$0.004 & 67.510$\pm$0.631 \\
EA-RG-MAPPO-S & 6 & 0.919$\pm$0.012 & 0.064$\pm$0.006 & 0.018$\pm$0.006 & 64.286$\pm$1.610 \\
EA-RG-MAPPO-S & 8 & 0.890$\pm$0.021 & 0.083$\pm$0.012 & 0.027$\pm$0.014 & 66.790$\pm$4.714 \\
EA-RG-MAPPO-S & 10 & 0.879$\pm$0.017 & 0.086$\pm$0.020 & 0.039$\pm$0.011 & 69.144$\pm$1.739 \\
\bottomrule
\end{tabular}
\end{table*}


Table~\ref{tab:final_comm_300_results} and Fig.~\ref{fig:final_success_en}--Fig.~\ref{fig:final_collision_en} report the final 300-episode evaluation. EA-RG-MAPPO-S achieves success rates from 0.879 to 0.926 across four communication radii and keeps collision rates between 0.054 and 0.086. At radius 4, MAPPO has a collision rate of $0.228\pm0.099$, while EA-RG-MAPPO-S reduces it to $0.054\pm0.007$. At radii 6 and 8, EA-RG-MAPPO-S obtains success/collision pairs of $0.919/0.064$ and $0.890/0.083$, respectively. Compared with GAT-MAPPO, EA-RG-MAPPO-S achieves higher success rates and lower collision rates at radii 8 and 10, with smaller standard deviations. These results indicate that standard graph attention alone does not guarantee cross-radius stability, while edge-aware attention and random-radius fine-tuning jointly provide a more robust coordination representation.

\begin{figure}[t]
    \centering
    \includegraphics[width=0.48\textwidth]{final_300_success_rate.png}
    \caption{Success rates under different communication radii in the 300-episode final evaluation.}
    \label{fig:final_success_en}
\end{figure}

\begin{figure}[t]
    \centering
    \includegraphics[width=0.48\textwidth]{final_300_collision_rate.png}
    \caption{Collision rates under different communication radii in the 300-episode final evaluation.}
    \label{fig:final_collision_en}
\end{figure}

\subsection{Ablation Analysis}

\begin{table*}[t]
\centering
\caption{Ablation study of role graph, edge-aware attention, and staged random-radius fine-tuning. Results are reported as mean$\pm$std over three seeds.}
\label{tab:ablation_results}
\begin{tabular}{lccccc}
\toprule
Method & Radius & Success $\uparrow$ & Collision $\downarrow$ & Timeout $\downarrow$ & Avg. steps $\downarrow$ \\
\midrule
RG-MAPPO & 4 & 0.907$\pm$0.042 & 0.080$\pm$0.044 & 0.017$\pm$0.006 & 73.200$\pm$7.600 \\
RG-MAPPO & 6 & 0.917$\pm$0.076 & 0.063$\pm$0.059 & 0.023$\pm$0.023 & 72.770$\pm$11.410 \\
RG-MAPPO & 8 & 0.887$\pm$0.146 & 0.097$\pm$0.125 & 0.017$\pm$0.021 & 72.000$\pm$13.450 \\
RG-MAPPO & 10 & 0.907$\pm$0.071 & 0.073$\pm$0.059 & 0.020$\pm$0.017 & 72.040$\pm$10.180 \\
\midrule
EA-RG-MAPPO & 4 & 0.927$\pm$0.021 & 0.037$\pm$0.015 & 0.037$\pm$0.006 & 65.850$\pm$2.780 \\
EA-RG-MAPPO & 6 & 0.887$\pm$0.015 & 0.073$\pm$0.032 & 0.040$\pm$0.036 & 63.410$\pm$2.200 \\
EA-RG-MAPPO & 8 & 0.900$\pm$0.036 & 0.083$\pm$0.012 & 0.020$\pm$0.026 & 63.780$\pm$4.580 \\
EA-RG-MAPPO & 10 & 0.837$\pm$0.032 & 0.113$\pm$0.015 & 0.050$\pm$0.017 & 73.360$\pm$4.120 \\
\midrule
EA-RG-MAPPO-S & 4 & 0.907$\pm$0.012 & 0.067$\pm$0.012 & 0.027$\pm$0.012 & 69.470$\pm$2.300 \\
EA-RG-MAPPO-S & 6 & 0.907$\pm$0.015 & 0.073$\pm$0.021 & 0.020$\pm$0.010 & 62.680$\pm$6.510 \\
EA-RG-MAPPO-S & 8 & 0.883$\pm$0.051 & 0.083$\pm$0.031 & 0.033$\pm$0.032 & 65.210$\pm$7.700 \\
EA-RG-MAPPO-S & 10 & 0.880$\pm$0.020 & 0.090$\pm$0.026 & 0.033$\pm$0.021 & 70.350$\pm$3.540 \\
\bottomrule
\end{tabular}
\end{table*}


Table~\ref{tab:ablation_results} reports the ablation results with 100 episodes per seed and is used for module analysis. The comparison among RG-MAPPO, EA-RG-MAPPO, and EA-RG-MAPPO-S indicates that relative edge features reduce collision under small communication radii and alleviate instability at radius 8 for some seeds. Staged random-radius fine-tuning improves generalization at radius 10 and makes the policy more balanced across multiple communication radii. This table has a different evaluation budget from the final 300-episode main table, so it is used as module-level evidence.

\subsection{Visualization Analysis}

Per-seed scatter plots show that MAPPO has larger seed-to-seed variation, while EA-RG-MAPPO-S has more concentrated success and collision values. Trajectory case studies show that baseline methods may produce inter-UAV collisions under the same environment seed, whereas EA-RG-MAPPO-S can maintain successful pursuit. Attention heatmaps show that the graph attention distribution changes with communication radius, supporting the interpretation that the policy adapts its information aggregation under limited communication.

\subsection{Target-Intent Branch Diagnostic}

The implementation previously included an auxiliary target-intent prediction branch. Diagnostic results show a plain accuracy of 0.587 and a balanced accuracy of 0.200, indicating that the branch mainly predicts the majority class. Therefore, this branch cannot support a high-accuracy intent-recognition claim and is not used as a main contribution. Future work on intent modeling should introduce short history, target turning-rate features, more balanced target maneuver sampling, and balanced accuracy as a primary metric.
